% GECo-CF working manuscript. Replace all TODO markers before submission.
\documentclass[runningheads]{llncs}

\usepackage[T1]{fontenc}
\usepackage{graphicx}
\usepackage{booktabs}
\usepackage{multirow}
\usepackage{amsmath,amssymb}
\usepackage{url}

\newcommand{\method}{GECo-CF}

\begin{document}

\title{GECo-CF: Community-Guided Coarse-to-Fine Counterfactual
Explanations for Graph Neural Networks}
\titlerunning{Community-Guided Counterfactual Explanations for GNNs}

% TODO: Replace the anonymous author block after double-blind review.
\author{Anonymous Authors}
\authorrunning{Anonymous Authors}
\institute{Anonymous Institution}

\maketitle

\begin{abstract}
Counterfactual explanations describe small graph changes that alter a trained
graph neural network's (GNN's) prediction. Existing methods commonly
search over the full computation graph or optimize a dense edge mask, which
can be inefficient and may fragment coherent graph structure. We propose
\method, a post-hoc, model-agnostic framework that uses communities as a
coarse structural prior for counterfactual search. Given a graph or a target
node, \method first measures the prediction sensitivity to interventions on
each detected community. It then expands the edge search space adaptively,
prioritizing boundary edges of influential communities before admitting
intra-community and global edges. A fine-grained, iteratively rescored search
finds a prediction-flipping edge set, after which reverse pruning removes
redundant edits. This design aims to retain the validity of global search while
reducing model queries and limiting damage to community structure. We first
establish a coarse-localization guarantee under bounded edge-interaction
residuals, then test its practical premise on synthetic BA-2Motifs and the
real-world Mutagenicity dataset. Across three community-detection and control
seeds, the top-ranked community covers a globally Top-5 sensitive edge in
$78.6$--$81.6\%$ and $85.3$--$87.3\%$ of eligible graphs, respectively,
compared with $52.1$--$53.2\%$ and $57.3$--$59.0\%$ for random selection.
Paired bootstrap intervals confirm the margin-drop and Top-5-coverage
advantages on both datasets. These results support community intervention as
a coarse localization signal; validity, minimality, and efficiency of the
complete search remain to be established.

\keywords{Graph neural networks \and Counterfactual explanations \and
Explainable AI \and Community structure}
\end{abstract}

\section{Introduction}
\label{sec:introduction}

Graph neural networks (GNNs) have become a standard tool for learning from
relational data, including molecules, social networks, and knowledge graphs
\cite{kipf2017gcn,xu2019gin}. Their predictions, however, are produced through
multi-hop message passing and are difficult to inspect directly. This lack of
transparency is consequential in scientific and high-stakes applications,
where a user needs not only an accurate prediction but also a concrete account
of which relational evidence supports it. Post-hoc GNN explainers address this
need by identifying influential subgraphs, features, or training examples
\cite{ying2019gnnexplainer,luo2020pgexplainer,yuan2022survey}. Most early
methods provide factual
explanations: they describe graph components that support the prediction that
was made.

Counterfactual explanations answer a different and actionable question:
\emph{what is the smallest plausible change to the input graph that would
change the model's prediction?} For a fixed predictor, a counterfactual graph
should flip the original label while remaining close to the observed graph.
Methods such as CF-GNNExplainer formulate this objective as edge-mask
optimization \cite{lucic2022cfgnn}, while subsequent approaches pursue robust
decision regions or use latent generative models and discrete diffusion to
improve plausibility and distributional alignment
\cite{bajaj2021rcexplainer,ma2022clear,huang2023d4explainer}. Despite this
progress, finding a high-quality structural counterfactual remains a difficult
combinatorial problem. A computation graph with $|E|$ editable edges already
induces $2^{|E|}$ deletion subsets, and unrestricted edit models are larger
still. Dense relaxation reduces the immediate optimization burden but does not
by itself identify where consequential edits should be sought.

Graphs are not merely collections of interchangeable edges. Many exhibit
mesoscopic organization in which nodes form densely connected communities
linked by comparatively sparse boundary edges \cite{fortunato2010community}.
Recent work has started to exploit communities or learned clusters for factual
graph explanations \cite{amato2026geco,yang2026ciflow}, suggesting that this
organization can provide a useful prior. Its role in counterfactual search,
however, remains underexplored. Intervening on a community can cheaply test
whether a coherent region is relevant to the current prediction; its boundary
can then provide a compact set of candidate edits. Yet naively restricting
edits to a single detected community is unsafe:
community detection may be unstable, predictive evidence may span several
communities, and the decisive edge may lie inside rather than between them.
The central challenge is therefore to exploit community structure without
sacrificing counterfactual validity.

We introduce \method, a community-guided coarse-to-fine framework for
structural counterfactual explanations. At the coarse stage, \method partitions
the explanation graph and ranks communities by the decrease in the original
class margin caused by a temporary community intervention. The score is
normalized for intervention size to reduce the systematic preference for
large communities. At the fine stage, the method first searches boundary edges
associated with highly ranked communities, then adaptively admits their
internal edges, additional communities, and finally the full graph when
necessary. Candidate edges are rescored after each edit to account for their
interactions. Once the prediction flips, reverse pruning restores every edge
whose deletion is unnecessary. Thus, communities guide the order of search
rather than impose a hard constraint on the returned explanation.

Our study is organized around three questions. First, do high-scoring
communities contain prediction-sensitive edges more often than random or
size-based choices? Second, can community-guided expansion reduce model queries
without materially lowering counterfactual validity? Third, does prioritizing
community boundaries yield explanations that preserve more of the original
mesoscopic structure? These questions separate the proposed mechanism from
the performance of the complete system and make a negative result
interpretable.

The main contributions are:
\begin{itemize}
  \item We formulate size-normalized community intervention as a coarse signal
  for locating counterfactually sensitive regions and prove, under bounded
  edge-interaction residuals, a lower bound on the most sensitive edge within
  a scored community.
  \item We develop an adaptive coarse-to-fine search that prioritizes community
  boundaries but falls back to broader candidate sets to preserve validity.
  \item We introduce an evaluation protocol that reports prediction flips and
  edit size together with query cost and community-structure damage across
  graph- and node-classification tasks.
  \item We empirically test the community-guidance hypothesis against random,
  largest-community controls, with a single-edge sensitivity oracle for
  evaluation. On BA-2Motifs and
  Mutagenicity, the top-ranked community covers a Top-5 sensitive edge in
  $78.6$--$81.6\%$ and $85.3$--$87.3\%$ of eligible graphs, versus
  $52.1$--$53.2\%$ and $57.3$--$59.0\%$ for random selection.
\end{itemize}

\section{Related Work}
\label{sec:related-work}

\paragraph{Factual GNN explanations.}
GNNExplainer learns an instance-specific mask over edges and features by
maximizing mutual information with the original prediction
\cite{ying2019gnnexplainer}. PGExplainer instead parameterizes mask generation
so that explanation knowledge can be shared across instances
\cite{luo2020pgexplainer}. Surveys and comparative studies organize the wider
design space and show that explainer behavior depends on the task, predictor,
and evaluation protocol \cite{yuan2022survey,longa2025explaining}. These methods
primarily identify evidence supporting an existing prediction; they do not
require the edited graph to produce a different prediction.

\paragraph{Counterfactual graph explanations.}
CF-GNNExplainer directly optimizes edge deletions for node-classification
counterfactuals \cite{lucic2022cfgnn}. RCExplainer seeks robust explanations
associated with shared decision regions \cite{bajaj2021rcexplainer}, whereas
CLEAR uses a graph variational autoencoder and D4Explainer uses discrete
denoising diffusion to improve plausibility or distributional alignment
\cite{ma2022clear,huang2023d4explainer}. These approaches motivate our common
requirements of validity and proximity. In contrast, our focus is the search
process: communities prioritize discrete candidates, while adaptive expansion
retains access to the complete editable edge set.

\paragraph{Community-aware graph explanations.}
GECo evaluates detected communities and combines the relevant ones into a
factual explanation for graph classification \cite{amato2026geco}. CIFlow-GNN
is an intrinsically interpretable architecture that controls information flow
between learned clusters \cite{yang2026ciflow}. Community structure has also
been incorporated into explanation objectives for fair heterogeneous link
prediction \cite{wen2025fexpcom}. These studies establish that mesoscopic
structure can organize graph explanations, but they do not address minimal
prediction-flipping edits. \method differs by remaining post-hoc and
model-agnostic, using community interventions only to order a counterfactual
search, and falling back to broader candidates when the prior is insufficient.

\section{Preliminaries and Problem Formulation}
\label{sec:preliminaries}

Let $G=(V,E,X)$ be an attributed graph and let $f_\theta$ be a trained,
frozen GNN. For graph classification, the explanation target $\tau$ is $G$;
for node classification, $\tau$ is a node $v$ and the explanation graph is its
$L$-hop computation graph for an $L$-layer message-passing GNN. Let
$y_0=\arg\max_c z_c(\tau;G)$ be the original predicted class, where $z_c$
denotes the logit for class $c$. We define the original-class margin as
\begin{equation}
  m_\tau(G)=z_{y_0}(\tau;G)-\max_{c\ne y_0}z_c(\tau;G).
\end{equation}
Under the deletion-only setting used in our primary evaluation, the objective
is to find a small set $E_{\mathrm{edit}}\subseteq E$ such that
$\arg\max_c f_\theta(G\setminus E_{\mathrm{edit}})_c\ne y_0$. Unless stated
otherwise, explanations are generated only for correctly classified test
instances and all explainers receive the same instances and edit budget.

\section{Method}
\label{sec:method}

\subsection{Community Intervention Score}
Let $\mathcal{C}=\{C_1,\ldots,C_K\}$ be a partition returned by Louvain
community detection \cite{blondel2008louvain}. For community $C_i$, let
$A(C_i)=E(C_i)\cup\delta(C_i)$ contain its internal and boundary edges. We
temporarily mask $A(C_i)$ and define
\begin{equation}
  \widetilde{s}_i=
  \frac{m_\tau(G)-m_\tau(G\setminus A(C_i))}{|A(C_i)|^\alpha},
  \qquad \alpha\geq 0.
\end{equation}
The intervention is used only to rank communities; it is not returned as the
counterfactual. We fix $\alpha=0.5$ for the Stage-1 hypothesis test; any later
tuning for the complete method is restricted to validation data.

\subsection{Coarse-Localization Guarantee}
For an editable edge $e$, define its individual margin sensitivity as
\begin{equation}
  d_e=m_\tau(G)-m_\tau(G\setminus\{e\}).
\end{equation}
For a community $C$, write its unnormalized intervention effect as
$D_C=m_\tau(G)-m_\tau(G\setminus A(C))$ and define the edge-interaction
residual
\begin{equation}
  r_C=D_C-\sum_{e\in A(C)}d_e.
\end{equation}
The residual captures all non-additive interactions introduced by deleting
the edges jointly. The following proposition formalizes the coarse-grained
role of the community score.

\begin{proposition}[Sensitive-edge lower bound]
\label{prop:localization}
If $|r_C|\leq\varepsilon_C$ for a nonempty intervention set $A(C)$, then
\begin{equation}
  \max_{e\in A(C)} d_e
  \geq |A(C)|^{\alpha-1}\widetilde{s}_C
  -\frac{\varepsilon_C}{|A(C)|}.
  \label{eq:sensitive-lower-bound}
\end{equation}
\end{proposition}

\begin{proof}
Because a maximum is at least the arithmetic mean,
$\max_{e\in A(C)}d_e\geq |A(C)|^{-1}\sum_{e\in A(C)}d_e$. By the residual
definition and $r_C\leq\varepsilon_C$, the sum is at least
$D_C-\varepsilon_C$. Finally,
$D_C=|A(C)|^\alpha\widetilde{s}_C$, which gives
Eq.~\eqref{eq:sensitive-lower-bound}.
\end{proof}

Proposition~\ref{prop:localization} does not assert an unconditional monotone
correlation for arbitrary nonlinear GNNs: large interaction residuals can
produce synergy or cancellation. Instead, it states why community
intervention is a valid coarse explanatory view when joint edge effects are
approximately additive or their residual is controlled. For communities of
equal intervention size and residual bound, ranking by $\widetilde{s}_C$
directly ranks a lower bound on the strongest contained edge. With unequal
sizes, $\alpha$ controls the trade-off between total intervention effect and
local sensitivity. Our experiments therefore test two empirical implications:
whether community scores correlate positively with contained edge sensitivity
and whether the top-ranked community covers globally sensitive edges more
often than size-based or random selection.

\subsection{Adaptive Fine-Grained Search}
Communities are ordered by $\widetilde{s}_i$. The candidate set initially
contains boundary edges of the top-ranked community. If no valid
counterfactual is found within the current stage, it is expanded to include
that community's internal edges, edges associated with subsequent communities,
and finally all editable edges. Within a stage, each candidate edge $e$ is
scored by the margin reduction caused by deleting it from the current graph.
The best edge is removed and the remaining edges are rescored until the label
changes or the budget is exhausted. Reverse pruning then adds deleted edges
back whenever doing so preserves the counterfactual label.

\section{Experimental Setup}
\label{sec:experiments}

\subsection{Datasets}
We begin with a deliberately small validation suite and expand only after the
community-guidance hypothesis is supported. Table~\ref{tab:datasets} lists the
planned datasets. All statistics will be computed from the exact processed
splits used in our code rather than copied across incompatible dataset
versions.

\begin{table}[t]
\centering
\caption{Datasets and their role. ``Completed'' denotes the Stage-1
community-localization evaluation reported in this paper; the remaining rows
describe the planned full-method evaluation.}
\label{tab:datasets}
\small
\begin{tabular}{llll}
\toprule
Task & Dataset & Source/type & Role \\
\midrule
Graph & BA-2Motifs & synthetic & Completed Stage-1 test \\
Graph & Mutagenicity & molecular & Completed Stage-1 test \\
Graph & MUTAG & molecular & Planned real-world test \\
Graph & Benzene & molecular & Real-world extension \\
\midrule
Node & Tree-Cycles & synthetic & Simple motif test \\
Node & Tree-Grid & synthetic & Complex motif test \\
Node & BA-Shapes & synthetic & Large local search space \\
Node & BA-Community & synthetic & Explicit community structure \\
Node & Cora/Citeseer & citation & Optional real-world extension \\
\bottomrule
\end{tabular}
\end{table}

\paragraph{Graph classification.}
BA-2Motifs is the primary diagnostic dataset because its class-dependent
motifs permit qualitative checks of whether highly ranked communities overlap
predictive structure. Mutagenicity tests whether the same localization signal
transfers to real molecular graphs without synthetic motif ground truth.
MUTAG and Benzene remain planned extensions for comparison with full
counterfactual explainers. Synthetic ground truth is used only for diagnostic
localization metrics; counterfactual quality is always evaluated against the
frozen predictor rather than assumed causal ground truth.

\paragraph{Node classification.}
Tree-Cycles, Tree-Grid, BA-Shapes, and BA-Community follow the family of
synthetic benchmarks used in prior GNN explanation work
\cite{ying2019gnnexplainer,lucic2022cfgnn}. BA-Community is especially relevant
for testing the proposed prior, whereas BA-Shapes stresses candidate-space
reduction in larger computation graphs. For each target node, we explain the
$L$-hop receptive field of the predictor. Communities are detected once on the
full graph and mapped to each local computation graph to avoid inconsistent
per-instance partitions. Cora and Citeseer are optional extensions; because
they lack reliable counterfactual ground truth, they are evaluated using
validity, minimality, efficiency, and structural preservation rather than
motif overlap.

\paragraph{Preprocessing and splits.}
We will report the number of graphs or nodes, class balance, feature
dimensionality, mean nodes and edges, split protocol, and the fraction of
eligible test instances for every processed dataset. Before counterfactual
evaluation, we will additionally report the mean number of communities,
community-size distribution, modularity, boundary-edge ratio, and the fraction
of explanation graphs containing only one community. Datasets with a dominant
single-community regime cannot provide meaningful evidence for the proposed
mechanism and will be identified as such rather than silently pooled.

\subsection{Stage-1 Community-Localization Pilot}
The first-stage pilot isolates the premise on which the subsequent
coarse-to-fine search depends: whether an inexpensive community-level
intervention localizes edges to which the frozen predictor is sensitive.
Table~\ref{tab:stage1-data} reports the exact processed datasets and frozen
predictors. The PyTorch Geometric version of BA-2Motifs contains 1,000 graphs,
10 node features, and two balanced classes. Its three-layer GCN uses 20 hidden
units per layer, raw adjacency aggregation (no normalization or added
self-loops), and concatenated global max and sum readouts. Mutagenicity
contains 4,337 molecular graphs with 14-dimensional categorical node features
and class counts 2,401/1,936. Its predictor is a three-layer GIN with 64 hidden
units, batch normalization, sum pooling, and dropout 0.5. Each dataset uses a
stratified 80/10/10 split and predictor seed 0. The frozen predictors reach
0.980 and 0.796 test accuracy, leaving 98 and 347 correctly classified test
graphs eligible for the pilot.

\begin{table}[t]
\centering
\caption{Processed Stage-1 datasets and frozen predictors. Eligible graphs are
correctly classified test instances.}
\label{tab:stage1-data}
\small
\begin{tabular}{lrrrrr}
\toprule
Dataset & Graphs & Features & Val. acc. & Test acc. & Eligible \\
\midrule
BA-2Motifs & 1,000 & 10 & 1.000 & 0.980 & 98 \\
Mutagenicity & 4,337 & 14 & 0.822 & 0.796 & 347 \\
\bottomrule
\end{tabular}
\end{table}

For every eligible graph, Louvain communities are ranked by the intervention
score with $\alpha=0.5$. We compare the top-ranked community with the largest
community and with a uniformly selected random community. The random control
is repeated 30 times per graph. We run pilot seeds 0, 1, and 2; these seeds
change Louvain tie-breaking and the random control but reuse the same seed-0
predictor and test split. They therefore measure robustness of the
localization procedure, not variation across independently trained GNNs.

As an evaluation oracle, we delete every edge individually and measure the
resulting decrease in the original-class margin. This exhaustive scan is not
part of the proposed search algorithm. It supplies the globally most
sensitive edge, the five most sensitive edges, and an edge-level ranking
against which community scores can be assessed. We report community
intervention margin drop, prediction-flip rate under the whole-community
intervention, coverage of the global Top-5 sensitive edges, and Spearman rank
correlation between community scores and aggregated edge sensitivity. Paired
differences over the same eligible graphs within each dataset are accompanied
by percentile bootstrap 95\% confidence intervals with 5,000 resamples.
Results are not pooled across datasets because their predictor margin scales
differ.

\subsection{Evaluation Protocol}
The primary metrics are validity (flip rate), number of edits, normalized edit
distance, and validity as a function of edit budget. Efficiency is measured by
end-to-end runtime, number of GNN forward passes, and number of candidate edges
examined. To assess structural preservation, we report candidate-space
reduction, the intra-community damage rate, the fraction of edits on boundary
edges, change in modularity under the original fixed partition, the number of
communities touched, and connectivity changes. Edit-size averages are reported
with validity and are not restricted to easy successful instances without an
explicit failure penalty.

We compare against Random-CF and unrestricted Global-Greedy-CF as diagnostic
controls, CF-GNNExplainer as an optimization-based baseline
\cite{lucic2022cfgnn}, and RCExplainer, CLEAR, and D4Explainer where their
task and edit settings permit a fair comparison
\cite{bajaj2021rcexplainer,ma2022clear,huang2023d4explainer}.
All methods use the same trained predictor, test instances, permitted edit
operations, and maximum budget. In Stage 1, three pilot seeds vary Louvain
tie-breaking and the random control while holding predictor seed 0 fixed;
instance-level comparisons use paired bootstrap confidence intervals. The
complete-method evaluation will additionally aggregate over independently
trained predictor seeds.

\section{Results}
\label{sec:results}

\subsection{Stage-1 Community Localization}
The graph partitions are non-degenerate and stable across the three pilot
seeds. On BA-2Motifs, the mean number of communities ranges from 4.92 to 4.97,
mean modularity from 0.6062 to 0.6067, and the mean cross-community edge ratio
from 0.1523 to 0.1543. The corresponding Mutagenicity ranges are 5.80--5.81,
0.5788--0.5793, and 0.2049--0.2059. No eligible graph in either dataset is
assigned to a single community. Thus, the pilot meaningfully tests a
community prior rather than a trivial single-block partition.

Table~\ref{tab:pilot-summary} summarizes the localization outcomes. On both
datasets, the intervention chosen by the proposed score produces a larger
margin drop and covers globally Top-5 sensitive edges more often than either
control. The mean community-score/edge-sensitivity Spearman correlation is
positive in every run: 0.447--0.467 on BA-2Motifs and 0.498--0.507 on
Mutagenicity. Together with Top-5 coverage of 0.803 and 0.863, respectively,
this supports the two empirical implications of
Proposition~\ref{prop:localization}: community scores provide useful
coarse-grained localization of fine-grained sensitivity.

\begin{table}[t]
\centering
\caption{Stage-1 results. Entries are means over pilot seeds 0, 1, and 2; the
same predictor and test split are reused within each dataset. Margin drops are
not compared across datasets because their logit scales differ.}
\label{tab:pilot-summary}
\small
\begin{tabular}{llrrr}
\toprule
Dataset & Metric & Top-ranked & Largest & Random \\
\midrule
\multirow{3}{*}{BA-2Motifs}
 & Margin drop & 4.885 & 1.184 & 1.290 \\
 & Whole-community flip rate & 0.500 & 0.480 & 0.350 \\
 & Global Top-5 edge coverage & 0.803 & 0.639 & 0.527 \\
\addlinespace
\multirow{3}{*}{Mutagenicity}
 & Margin drop & 1.911 & 0.956 & 0.721 \\
 & Whole-community flip rate & 0.359 & 0.162 & 0.112 \\
 & Global Top-5 edge coverage & 0.863 & 0.748 & 0.582 \\
\bottomrule
\end{tabular}
\end{table}

The paired analysis in Table~\ref{tab:pilot-bootstrap} confirms that the
advantages do not arise from comparing different graph sets. On BA-2Motifs,
the margin-drop and Top-5-coverage intervals exclude zero for every pilot seed,
as does the flip-rate comparison with random selection. Its flip-rate
advantage over the largest community remains only 0.020, with an interval
including zero. On Mutagenicity, all six paired effects are positive and every
per-seed 95\% interval excludes zero. In particular, top-ranked selection
improves margin drop by 0.917--0.976 and Top-5 coverage by 0.098--0.133 over
the largest-community control. The real-data replication therefore rules out
the simple explanation that the method merely selects large communities.

\begin{table}[t]
\centering
\caption{Paired effects of top-ranked community selection. The mean is the
range across pilot seeds; the CI envelope is the smallest interval containing
the three per-seed 95\% percentile bootstrap intervals.}
\label{tab:pilot-bootstrap}
\small
\begin{tabular}{llrr}
\toprule
Dataset & Effect & Mean range & 95\% CI envelope \\
\midrule
\multirow{4}{*}{BA-2Motifs}
 & Margin vs. largest & 3.609--3.851 & [2.918, 4.520] \\
 & Margin vs. random & 3.550--3.641 & [3.208, 3.975] \\
 & Top-5 vs. largest & 0.153--0.184 & [0.020, 0.316] \\
 & Top-5 vs. random & 0.264--0.290 & [0.173, 0.377] \\
\addlinespace
\multirow{6}{*}{Mutagenicity}
 & Margin vs. largest & 0.917--0.976 & [0.793, 1.111] \\
 & Margin vs. random & 1.172--1.203 & [1.083, 1.297] \\
 & Flip vs. largest & 0.190--0.207 & [0.150, 0.251] \\
 & Flip vs. random & 0.242--0.257 & [0.207, 0.293] \\
 & Top-5 vs. largest & 0.098--0.133 & [0.046, 0.184] \\
 & Top-5 vs. random & 0.263--0.290 & [0.223, 0.327] \\
\bottomrule
\end{tabular}
\end{table}

\subsection{Discussion}
The theoretical and empirical results answer complementary questions.
Proposition~\ref{prop:localization} establishes a conditional guarantee: when
joint deletion effects are sufficiently close to the sum of their individual
effects, a high community score certifies a lower bound on the sensitivity of
at least one contained edge. It explains why community intervention can serve
as a coarse explanatory lens, but does not guarantee a positive rank
correlation for every nonlinear GNN. The experiments evaluate whether this
conditional mechanism remains informative when its interaction residuals are
unknown. Positive Spearman correlations on both a controlled synthetic
dataset and real molecular graphs, together with consistently improved Top-5
coverage, provide convergent evidence that it does.

The largest-community control is important for interpreting this evidence.
Because $D_C$ can grow with the number of intervened edges, an unnormalized
method could appear effective merely by selecting a large region. The
$|A(C)|^{0.5}$ normalization instead selects communities whose margin effect
is large relative to intervention size. Its advantage over the largest
community in margin reduction and sensitive-edge coverage on both datasets,
and additionally in whole-community flip rate on Mutagenicity, shows that the
signal is not reducible to community size. Stability across three Louvain and
random-control seeds further suggests that the result is not driven by one
partition tie-break. These findings justify using the ranking to prioritize a
fine edge search; they do not yet establish that this search is faster or
returns smaller counterfactuals.

\paragraph{Scope and limitations.}
This pilot supports only the first research question: community intervention
can localize prediction-sensitive regions on one synthetic and one real-world
graph-classification dataset. It does not yet show that \method produces
valid, sparse, or structurally preserving counterfactual explanations.
Removing all edges associated with a community is a diagnostic perturbation
rather than the returned explanation, and the single-edge scan is an
exhaustive evaluation oracle rather than an efficient search. Each dataset
currently uses one trained predictor. The three pilot seeds assess
community-detection and random-control stability, not predictor-seed
robustness. Moreover, Proposition~\ref{prop:localization} depends on a bounded
interaction residual and does not claim universal positive correlation.
Evaluation of the adaptive fine search, multiple trained predictors,
additional datasets, full counterfactual baselines, and the proposed validity,
minimality, efficiency, and structural-preservation metrics is left to the
next experimental stage.

\section{Conclusion}
\label{sec:conclusion}
We presented \method, a community-guided coarse-to-fine framework for
counterfactual GNN explanations, established a conditional coarse-localization
guarantee, and evaluated its practical premise in a controlled first-stage
pilot. On BA-2Motifs and Mutagenicity, intervention-ranked communities
consistently yielded larger margin reductions and higher coverage of globally
sensitive edges than random and largest-community controls. Positive
community/edge sensitivity correlations on both datasets further support the
community score as a coarse explanatory signal. This justifies using
community interventions to order a subsequent edge-level search, but it is not
evidence that the complete method already returns minimal counterfactuals. The
next stage must evaluate the adaptive search itself across predictor seeds,
datasets, and counterfactual baselines before claims about validity,
efficiency, or structural preservation can be made.

\bibliographystyle{splncs04}
\bibliography{references}

\end{document}
